\section{Categorical fusion paths and regrouping}
\label{sec:intrinsic-categories}

In the categorical construction, the bond label carries a fusion-path
space. It cannot be absorbed into a classical flag by the group
substitution used above. Nevertheless, a fixed regrouping determines
exactly where its weights enter the canonical tensor.

\subsection{Base tensors and recoupling conventions}

The Ising rules are
\begin{align}
	\psi^2     & =1,                 &
	\psi\sigma & =\sigma\psi=\sigma, &
	\sigma^2   & =1+\psi,
\end{align}
with dimensions $d_1=d_\psi=1$, $d_\sigma=\sqrt2$.
The Fibonacci rules are
\begin{align}
	\tau^2 & =1+\tau,                  &
	d_1    & =1,                       &
	d_\tau & =\phi=\frac{1+\sqrt5}{2}.
\end{align}
Write $N_{ab}^c$ for the fusion coefficients and
$\mc D^2=\sum_ad_a^2$, equal to $4$ and $\phi+2$, respectively.

Use real associators $[F^{abc}_d]_{ef}$, zero unless
$N_{ab}^eN_{ec}^dN_{bc}^fN_{af}^d=1$. The nontrivial entries are
\begin{align}
	F^{\sigma\sigma\sigma}_\sigma
	 & =\frac1{\sqrt2}
	\begin{pmatrix}1&1\\1&-1\end{pmatrix},             \\
	[F^{\sigma\psi\sigma}_\psi]_{\sigma\sigma}
	 & =[F^{\psi\sigma\psi}_\sigma]_{\sigma\sigma}=-1, \\
	F^{\tau\tau\tau}_\tau
	 & =\begin{pmatrix}
		    \phi^{-1}   & \phi^{-1/2} \\
		    \phi^{-1/2} & -\phi^{-1}
	    \end{pmatrix}.
	\label{eq:F-tables}
\end{align}
All other admissible scalar entries are one. These are the usual
unitary fusion tables in the indicated gauge~\cite{Bultinck}.
Their orthogonality uses
$1/2+1/2=1$ and $\phi^{-2}+\phi^{-1}=1$.

Define the orthonormal path bases
\begin{align}
	\mc P_a & =\operatorname{span}\{\ket{x,y}:N_{ax}^y=1\},  \\
	\mc J   & =\{(x',\rho,x):N_{x'\rho}^x=1\},               \\
	A_a^{(y',\rho,y),(x',\rho,x)}
	        & =[F^{a x'\rho}_{y}]_{y'x}\ket{x',y'}\bra{x,y}.
	\label{eq:base-A}
\end{align}
Unequal external labels $\rho$ give zero letters. All bases use
lexicographic order with labels $(1,\psi,\sigma)$ or $(1,\tau)$.
The dimensions of $\mc P_a$ are $(3,3,4)$ and $(2,3)$.

On matching paths define
\begin{align}
	U_{ab}(\ket{z,y}_a\otimes\ket{x,z}_b)
	 & =\sum_c[F^{abx}_y]_{cz}\ket{c;x,y},
	\label{eq:path-U}
\end{align}
and set $U_{ab}=0$ on nonmatching paths.
Orthogonality gives $U_{ab}U_{ab}^\dagger=\Id$ and
$U_{ab}^\dagger U_{ab}=P_{ab}$, the matching-path projection.

The recoupling equation in these coordinates is
\begin{align}
	 & \sum_{z',z}
	   [F^{abx'}_{y'}]_{cz'}
	   [F^{a z'\rho}_{y}]_{y'z}
	   [F^{b x'\rho}_{z}]_{z'x}
	   [F^{abx}_{y}]_{dz}
	\nonumber                                      \\
	 & \qquad=\delta_{cd}[F^{c x'\rho}_{y}]_{y'x}.
	\label{eq:recoupling}
\end{align}
This is the pentagon change of parenthesization, with both output
labels $c,d$ retained. It can also be checked directly from
\cref{eq:F-tables}: use $t=\sqrt2$ and $t^2=2$ for Ising, and
$t=\phi^{-1/2}$ and $t^4+t^2=1$ for Fibonacci. The admissible domain is
\begin{align}
	N_{ab}^{c}N_{ab}^{d}N_{x'\rho}^{x}N_{y'\rho}^{y}
	N_{cx'}^{y'}N_{dx}^{y}=1.
\end{align}
There are $136$ Ising coordinates, including $12$ with $c\neq d$,
and $50$ Fibonacci coordinates, including $10$ with $c\neq d$.
Inserting the complete tables and reducing powers of $t$ makes
each residual zero. For example, the off-diagonal Fibonacci
coordinate with all six external nonunit labels equal to $\tau$
and $(c,d)=(1,\tau)$ is
\begin{align}
	-t^5+t^3-t^7
	 & =-t^5+t^3-(t^3-t^5)=0.
\end{align}
The zero-entry rule and the admissibility domain are essential;
no permutation symmetry of the F indices is assumed.

Expansion of the concatenated letters gives
\begin{align}
	(A_aA_b)^{hh'}
	 & =\sum_{z',z}
	    [F^{a z'\rho}_y]_{y'z}
	    [F^{b x'\rho}_z]_{z'x}
	\ket{z',y'}\ket{x',z'}
	\bra{z,y}\bra{x,z}.
\end{align}
Every ket and bra is matching. Applying \cref{eq:path-U} and
\cref{eq:recoupling} therefore proves the full-space identity
\begin{align}
	A_aA_b
	 & =U_{ab}^\dagger
	\left(\bigoplus_c\Id_{N_{ab}^c}\otimes A_c\right)U_{ab}.
	\label{eq:base-fusion}
\end{align}
In particular, the complement is zero before compression.

\subsection{All-length positivity and the weighted tensor}

Let $R(a_1,\ldots,a_k)$ denote the horizontal periodic contraction
of the indicated word. Direct substitution into \cref{eq:base-A}
gives the one-letter closures
\begin{align}
	R_1^{\mathrm I}
	                   & =\begin{pmatrix}1&1&1\\1&1&1\\1&1&2\end{pmatrix}, &
	R_\psi^{\mathrm I} & =R_\sigma^{\mathrm I}=0,                            \\
	R_1^{\mathrm F}
	                   & =\begin{pmatrix}1&1\\1&2\end{pmatrix},            &
	R_\tau^{\mathrm F}
	                   & =\phi^{-2}\ket{\tau,\tau}\bra{\tau,\tau}.
	\label{eq:seeds}
\end{align}
The nonzero Ising eigenvalues solve $(z-2)^2-2=0$ and are
$2\pm\sqrt2$. The Fibonacci unit-sector eigenvalues solve
$(z-1)(z-2)-1=0$ and are $\phi^2,\phi^{-2}$. All seeds are
positive.

Repeated application of \cref{eq:base-fusion} gives an isometry
$J_{\boldsymbol a}$ with
\begin{align}
	R(\boldsymbol a)
	 & =J_{\boldsymbol a}
	\left(\bigoplus_c\Id_{N_{\boldsymbol a}^c}\otimes R_c\right)
	J_{\boldsymbol a}^\dagger,            \\
	\tr R(\boldsymbol a)
	 & =\sum_cN_{\boldsymbol a}^c\tr R_c.
	\label{eq:word-spectrum}
\end{align}
Here $N_{\boldsymbol a}^c$ counts successive fusion paths. Since
the recoupling maps do not depend on the contracted horizontal
indices, they pass through their finite sums. This proves positivity
at every word length.

Let $\mc L=\mc R=\C^4$. Choose mutually orthogonal isometries
$\Phi_b:\mc P_b\to\mc R\otimes\mc L$ by assigning the ordered
basis vectors successively to
\begin{align}
	\ket{\varphi_{u,v}}
	 & =\frac12\sum_{r\in\Z_2^2}(-1)^{u\cdot r}\ket{r,r+v},
\end{align}
ordered first by $u$, then by $v$. Ten register vectors are used
for Ising and five for Fibonacci; the unused dimensions are six
and eleven.

For $\lambda_b>0$, put
\begin{align}
	\Lambda & =\sum_b\lambda_bd_b,         &
	w_f     & =\frac{d_f}{\mc D^2\Lambda}.
\end{align}
Writing $[Z]_{rs}=(\bra r\otimes\Id)Z(\ket s\otimes\Id)$,
the register-crossing tensor is
\begin{align}
	B_f^{(h,r,s),(h',p,q)}
	          & =\sum_{b,k}\lambda_b
	                       [\Phi_bA_b^{hk}\Phi_b^\dagger]_{rs}
	\otimes A_f^{kh'}\otimes\ket p\bra q,                      \\
	M^\lambda & =\bigoplus_fw_fB_f.
	\label{eq:crossing-category}
\end{align}
Its ring is a positive weighted sum of embedded words
$R(b_1,f_1,\ldots,b_N,f_N)$ by \cref{eq:word-spectrum}.

\subsection{Exact regrouping and canonical representatives}

Regroup the factors into $(R_{n-1},L_n,\mc P_{f_n})$.
The shifted tensor is
\begin{align}
	C^{hh'}
	 & =\bigoplus_f\sum_{b,k}w_f\lambda_b
	\Phi_bA_b^{hk}\Phi_b^\dagger\otimes A_f^{kh'}.
	\label{eq:shift-category}
\end{align}
The identity
$\sum_{r,s}\ket r\bra s\otimes[Z]_{rs}=Z$
proves equality of the periodic families at every $N$ under this
fixed regrouping.

Let $J_{bf}^{c}$ be the column of $U_{bf}^\dagger$ associated
with output $c$, and let
$V_{bf}^{c}=(\Phi_b\otimes\Id)J_{bf}^{c}$, placed in physical
sector $f$. Then
\begin{align}
	C^{hh'}
	 & =\sum_{b,f,c}w_f\lambda_b
	V_{bf}^{c}A_c^{hh'}(V_{bf}^{c})^\dagger,       \\
	\eta_c
	 & =\bigoplus_{b,f}w_f\lambda_b\Id_{N_{bf}^c}.
	\label{eq:category-eta}
\end{align}
The ranges are orthogonal. Their complement consists of unused
register states and nonmatching paths and is zero on every letter.

The canonical status of the $A_c$ must be checked independently of
this decomposition. Every nonzero letter in \cref{eq:base-A} is a
scalar matrix unit, so its transfer map kills off-diagonal matrices.
On diagonal matrices the transition matrix is
\begin{align}
	K_a[(x',y'),(x,y)]
	 & =\sum_\rho\left|[F^{a x'\rho}_y]_{y'x}\right|^2.
\end{align}
The full matrices are
\begin{align}
	K_1^{\mathrm I}=K_\psi^{\mathrm I}
	 & =\begin{pmatrix}1&1&1\\1&1&1\\1&1&2\end{pmatrix}, \\
	K_\sigma^{\mathrm I}
	 & =\begin{pmatrix}
		    1   & 1   & 1 & 1 \\
		    1   & 1   & 1 & 1 \\
		    1/2 & 1/2 & 1 & 1 \\
		    1/2 & 1/2 & 1 & 1
	    \end{pmatrix},                               \\
	K_1^{\mathrm F}
	 & =\begin{pmatrix}1&1\\1&2\end{pmatrix},            \\
	K_\tau^{\mathrm F}
	 & =\begin{pmatrix}
		    1         & 1 & 1           \\
		    \phi^{-2} & 1 & \phi^{-1}   \\
		    \phi^{-1} & 1 & 1+\phi^{-2}
	    \end{pmatrix}.
	\label{eq:base-transfer}
\end{align}
For example, the last entry is $1^2+(-\phi^{-1})^2$.
Every entry is positive, so every physical matrix unit occurs among
the letters. Each tuple is one-site injective and its transfer map
is primitive.

For Ising, vectors constant on the first two coordinates of $K_1$
and $K_\psi$ give the restriction
$\left(\begin{smallmatrix}2&1\\2&2\end{smallmatrix}\right)$.
For $\sigma$, vectors constant on the first and last pairs give
$\left(\begin{smallmatrix}2&2\\1&2\end{smallmatrix}\right)$.
Both characteristic polynomials are $z^2-4z+2$, with positive
Perron eigenvectors and root $2+\sqrt2$. For Fibonacci,
\begin{align}
	K_1^{\mathrm F}\binom1\phi
	 & =\binom{1+\phi}{1+2\phi}
	=\phi^2\binom1\phi,         \\
	K_\tau^{\mathrm F}
	\begin{pmatrix}\phi\\1\\ \phi\end{pmatrix}
	 & =\begin{pmatrix}
		    2\phi+1     \\
		    \phi^{-1}+2 \\
		    2+\phi+\phi^{-1}
	    \end{pmatrix}
	=\phi^2\begin{pmatrix}\phi\\1\\ \phi\end{pmatrix}.
\end{align}
Thus the common normalizations are
\begin{align}
	\kappa_{\mathrm I} & =\sqrt{2+\sqrt2}, &
	\kappa_{\mathrm F} & =\phi,            &
	\widehat A_c       & =A_c/\kappa.
	\label{eq:kappa}
\end{align}
No unexamined off-diagonal transfer modes remain.

The trace vectors distinguish all sectors. At
$h=(a,1,a)$ and $h'=(1,1,1)$,
\begin{align}
	\tr A_c^{hh'} & =\delta_{ac},
\end{align}
because admissibility requires $c1=a$ and the unit associator is
one. Hence no pair is related by similarity and a nonzero scalar.

Horizontally, the external $\rho$ blocks have dimensions $(3,3,4)$
or $(2,3)$. At fixed $\rho$, the nonzero F entries supply every
horizontal matrix unit: the physical ket and bra determine both
horizontal endpoints. The positive coefficients and orthogonal
embeddings in \cref{eq:category-eta} preserve that span. The blocks
are therefore injective. The equal-dimensional Ising blocks are
inequivalent because the physical trace vanishes for $\rho=\psi$
but not for $\rho=1$. The relevant cancellation is $1-1=0$.
For Fibonacci, the nonunit block has cancellation
$1-\phi\phi^{-1}=0$. All other pairs have different dimensions.
In particular, the shifted tensor is horizontally canonical and
nonsimple.

\begin{proposition}[Canonical data after categorical regrouping]
	\label{prop:category-shift}
	For every strictly positive weight vector, the shifted tensor has
	vertical BNT $\{\widehat A_c\}$ and data
	\begin{align}
		\mu_c^{\mathrm{sh}} & =\kappa\eta_c,              &
		m_c^{\mathrm{sh}}   & =\kappa\frac{d_c}{\mc D^2},   \\
		\chi_{abc}^{\mathrm{sh}}
		                    & =\kappa^{-1}\Id_{N_{ab}^c}, &
		c_{abc}^{\mathrm{sh}}(L)
		                    & =N_{ab}^c\kappa^{-L}.
		\label{eq:category-shift-data}
	\end{align}
	It generates positive trace-one rings and satisfies full-space RFP
	identities. The raw representatives $A_c$ have constant integer
	fusion coefficients.
\end{proposition}

\begin{proof}
	Only the masses, normalization, and channel construction remain.
	In the fusion ring, set $\mathsf d=\sum_ad_aa$. The explicit rules give
	$b\mathsf d=d_b\mathsf d$ and $\mathsf d^2=\mc D^2\mathsf d$.
	Consequently,
	\begin{align}
		\tr\eta_c
		 & =\frac1{\mc D^2\Lambda}
		\sum_{b,f}\lambda_bd_fN_{bf}^c \\
		 & =\frac1{\mc D^2\Lambda}
		\sum_b\lambda_bd_bd_c
		=\frac{d_c}{\mc D^2}.
		\label{eq:category-mass}
	\end{align}
	The blocked raw multiplicity is
	$\zeta_c=\bigoplus_{a,b}\eta_a\otimes\eta_b
		\otimes\Id_{N_{ab}^c}$, so
	\begin{align}
		\tr\zeta_c
		 & =\frac1{\mc D^4}\sum_{a,b}d_ad_bN_{ab}^c
		=\frac{d_c}{\mc D^2}
		=\tr\eta_c.
	\end{align}
	Its embeddings are compositions of $V_{bf}^c$ and
	$U_{ab}^\dagger$. The complete complements are the orthogonal
	complements of these ranges. \Cref{lem:channels} supplies the
	claimed CPTP maps, including their completion branches.

	For the trace, put $\Theta=\sum_b\lambda_bb$ and
	$W=\mathsf d/(\mc D^2\Lambda)$. Then
	\begin{align}
		\Theta W     & =\mathsf d/\mc D^2, &
		(\Theta W)^N & =\mathsf d/\mc D^2.
	\end{align}
	Using $\ell(c)=\tr R_c$ and \cref{eq:word-spectrum},
	\begin{align}
		\tr\rho_N(C) & =\ell(\mathsf d)/\mc D^2.
	\end{align}
	For Ising, $\ell(\mathsf d)=4$. For Fibonacci,
	$\ell(\mathsf d)=3+\phi\phi^{-2}=3+\phi^{-1}=\phi+2$.
	Both traces equal one. Finally, division of
	\cref{eq:base-fusion} by $\kappa^2$ gives
	\cref{eq:category-shift-data}.
\end{proof}

\subsection{Multiplicity spectra and the fixed-cut comparison}

The shifted raw multiplicities, including every copy, are
\begin{align}
	\eta_1^{\mathrm I}=\eta_\psi^{\mathrm I}
	 & =\frac1{4\Lambda}
	\diag(\lambda_1,\lambda_\psi,\sqrt2\lambda_\sigma), \\
	\eta_\sigma^{\mathrm I}
	 & =\frac1{4\Lambda}
	\diag(\sqrt2\lambda_1,\sqrt2\lambda_\psi,
	\lambda_\sigma,\lambda_\sigma),                     \\
	\eta_1^{\mathrm F}
	 & =\frac1{(\phi+2)\Lambda}
	\diag(\lambda_1,\phi\lambda_\tau),                  \\
	\eta_\tau^{\mathrm F}
	 & =\frac1{(\phi+2)\Lambda}
	\diag(\phi\lambda_1,\lambda_\tau,\phi\lambda_\tau).
	\label{eq:category-spectra}
\end{align}
For example, the three Fibonacci copies of $\tau$ come from
$(b,f)=(1,\tau),(\tau,1),(\tau,\tau)$. These spectra are not
identical across sectors.

For comparison, the original register-crossing construction has
\begin{align}
	B_fB_{f'}
	 & \simeq\bigoplus_{b,c}\lambda_b
	\Id_{N_{fbf'}^c}\otimes B_c,
	\label{eq:crossing-fusion}
\end{align}
where $N_{fbf'}^c=\sum_eN_{fb}^eN_{ef'}^c$.
The isometry first reads the inner register with $\Phi_b^\dagger$
and then applies $U_{fb}$ and $U_{ef'}$. Unused register vectors
and nonmatching paths give the zero complement.

At
\begin{align}
	\lambda_{\mathrm I} & =(1/2,3/10,1/5), &
	\lambda_{\mathrm F} & =(3/5,2/5),
	\label{eq:certified-weights}
\end{align}
the crossing sectors are canonical, with common transfer radii
\begin{align}
	\nu_{\mathrm I}^2 & =\frac{40+\sqrt{802}}{50}, &
	\nu_{\mathrm F}^2 & =\frac{47+5\sqrt{53}}{50}.
	\label{eq:crossing-radii}
\end{align}
The distinction between these radii and \cref{eq:kappa} is
essential: they normalize different tensors at different cuts.

Here is a compact specification of the finite canonical checks
behind this comparison. Define
\begin{align}
	X_f^{h,r,s,h'}
	 & =\sum_{b,k}\lambda_b
	              [\Phi_bA_b^{hk}\Phi_b^\dagger]_{rs}\otimes A_f^{kh'}, \\
	B_f^{(h,r,s),(h',p,q)}
	 & =X_f^{h,r,s,h'}\otimes\ket p\bra q.
\end{align}
For a basis index $i$ of $\mc P_b$, let $(u_{bi},v_{bi})$ be its
assigned Bell label and put
\begin{align}
	\beta_{bi}(r,l)
	 & =\tfrac12(-1)^{u_{bi}\cdot r}\delta_{l,r+v_{bi}}, \\
	[X_f^{h,r,s,h'}]_{(l,i),(m,j)}
	 & =\sum_{b,k,p,q}\lambda_b[A_b^{hk}]_{pq}
	\beta_{bp}(r,l)\beta_{bq}(s,m)[A_f^{kh'}]_{ij}.
	\label{eq:X-array}
\end{align}
This defines every entry. The ranks of the arrays of vectorized
$X_f$ are $(144,144,256)$ and $(64,144)$, their full column
dimensions. One nonzero maximal-minor certificate uses the field
reductions
\begin{align}
	\sqrt2      & \longmapsto327535\pmod{999983}, \\
	\phi^{-1/2} & \longmapsto441162\pmod{999809}.
\end{align}
Scanning lexicographic nonzero rows, retaining the first independent
rows and the first available pivot column at each step, gives
determinants
\begin{align}
	 & (266954,266954,254623)\pmod{999983}, \\
	 & (755529,823417)\pmod{999809}.
\end{align}
The defining field polynomials vanish at these images and all
denominators are invertible. Nonzero reduced determinants therefore
certify nonzero characteristic-zero minors, not merely numerical
ranks.

The horizontal fibers at fixed $\rho,h'$ are obtained by treating
$(f,i,j)$ as rows and $(h,r,s)$ as columns in \cref{eq:X-array}.
Their full ranks are $48,48,64$ for Ising sectors and $32,48$
for Fibonacci sectors. In lexicographic order of the global $h'$
indices, the corresponding determinant residues are
\begin{align}
	\mathrm I,\ \rho=1      & : (48068,48068,48068),\nonumber    \\
	\mathrm I,\ \rho=\psi   & : (731658,731658,731658),\nonumber \\
	\mathrm I,\ \rho=\sigma & : (125935,125935,741185,741185),   \\
	\mathrm F,\ \rho=1      & : (86291,86291),\nonumber          \\
	\mathrm F,\ \rho=\tau   & : (142240,916912,25178).
\end{align}
The same row-selection rule applies. Independent right matrix units
complete the horizontal and vertical spans. Trace-vector minors on
coordinates $(H,H')=(0,0),(0,48),(0,96)$ for Ising and
$(0,0),(0,32)$ for Fibonacci have residues $117748$ and $52490$,
respectively, where $H=16h+4r+s$. Thus the vertical sectors are
inequivalent.

The radii in \cref{eq:crossing-radii} follow from positive
two-dimensional invariant subspaces of the full transfer map.
The restriction matrices are
\begin{align}
	T_{1,\psi}^{\mathrm I}
	 & =\begin{pmatrix}19/25&21/50\\19/25&21/25\end{pmatrix}, &
	T_\sigma^{\mathrm I}
	 & =\begin{pmatrix}21/25&19/25\\21/50&19/25\end{pmatrix},   \\
	T_1^{\mathrm F}
	 & =\begin{pmatrix}13/25&17/25\\13/25&34/25\end{pmatrix}, &
	T_\tau^{\mathrm F}
	 & =\begin{pmatrix}34/25&13/25\\17/25&13/25\end{pmatrix}.
\end{align}
The subspaces consist of diagonal operators constant on path groups
$(\{0,1\},\{2\})$, $(\{0,1\},\{2,3\})$,
$(\{0\},\{1\})$, and $(\{0,2\},\{1\})$, respectively, tensored
with identities on both register halves. Substitution into
\cref{eq:X-array} verifies invariance on the full output, including
zero off-diagonal entries. The characteristic polynomials are
\begin{align}
	z^2-\frac85z+\frac{399}{1250},\qquad
	z^2-\frac{47}{25}z+\frac{221}{625}.
\end{align}
Their larger roots are \cref{eq:crossing-radii}. The restriction
Perron vectors give strictly positive definite eigenmatrices of
the full transfer maps. Conjugating by their square roots and
dividing by the eigenvalue makes each map unital, which proves
that these roots are the full spectral radii.

Thus at \cref{eq:certified-weights}, the crossing BNT data are
\begin{align}
	A_f^{\mathrm{cross}}   & =B_f/\nu,                                    &
	\mu_f^{\mathrm{cross}} & =\nu w_f,                                      \\
	\chi_{ff'c}^{\mathrm{cross}}
	                       & =\bigoplus_b(\lambda_b/\nu)\Id_{N_{fbf'}^c}.
\end{align}
The mass equation follows from
$\sum_{f,f'}d_fd_{f'}N_{fbf'}^c=\mc D^2d_bd_c$:
\begin{align}
	\sum_{f,f'}m_fm_{f'}\tr\chi_{ff'c}
	 & =\nu\sum_{f,f',b}w_fw_{f'}\lambda_bN_{fbf'}^c
	=\nu w_c.
\end{align}
In particular,
\begin{align}
	\chi_{\sigma\sigma1}^{\mathrm{cross}}
	 & =\nu^{-1}\diag(\lambda_1,\lambda_\psi),              \\
	\chi_{\tau\tau1}^{\mathrm{cross}}
	 & =\nu^{-1}\diag(\lambda_1,\lambda_\tau),              \\
	\chi_{\tau\tau\tau}^{\mathrm{cross}}
	 & =\nu^{-1}\diag(\lambda_1,\lambda_\tau,\lambda_\tau).
\end{align}
Their unequal eigenvalues prove the fixed-cut rescaling obstruction.
The canonical comparison is made only at the certified weights.
The all-positive-weight result in \cref{prop:category-shift} concerns
the shifted tensor and does not extend these crossing certificates
to every weight vector.

\subsection{What the categorical obstructions establish}

\begin{proposition}[Restricted spectator obstruction]
	\label{prop:restricted-obstruction}
	For strictly positive weights, the shifted categorical multiplicities
	cannot have a common mixed tensor factor under a sector-preserving
	one-site unitary. The periodic states are also not products of
	identical independent bond states, even after fixed blocking and
	unitary changes of coordinates.
\end{proposition}

\begin{proof}
	Suppose the first factorization were
	\begin{align}
		\eta_c & \simeq\sigma\otimes\xi_c
	\end{align}
	for every sector, allowing arbitrary residual factors $\xi_c$.
	Then $\rank\sigma$ divides every $\rank\eta_c$.
	The ranks in \cref{eq:category-spectra} are $3,3,4$ for Ising
	and $2,3$ for Fibonacci. Their greatest common divisor is one.
	Thus $\rank\sigma=1$, excluding a common mixed factor.

	For the second statement, positivity of the weights makes the ring
	ranks independent of their values. The multiplicity-counting element
	for Ising is $H=3+3\psi+4\sigma$. Only $R_1$ has nonzero rank,
	equal to two. Evaluation at the three fusion-ring characters gives,
	for $N\geq1$,
	\begin{align}
		\rank\rho_N^{\mathrm I}
		 & =2[1]H^N                                \\
		 & =\frac{(6+4\sqrt2)^N+(6-4\sqrt2)^N}{2}.
		\label{eq:Ising-rank}
	\end{align}
	Indeed, the two characters with $\psi=1$ have
	$\sigma=\pm\sqrt2$, and the character with $\psi=-1$ has
	$\sigma=0$ and $H=0$.

	For Fibonacci, write
	$(2+3\tau)^N=p_N+q_N\tau$. Its two character evaluations are
	\begin{align}
		p_N+\phi q_N     & =(2+3\phi)^N,      \\
		p_N-\phi^{-1}q_N & =(2-3\phi^{-1})^N.
	\end{align}
	Since $\rank R_1=2$, $\rank R_\tau=1$, and
	$\phi-\phi^{-1}=1$,
	\begin{align}
		\rank\rho_N^{\mathrm F}
		 & =2p_N+q_N                      \\
		 & =(2+3\phi)^N+(2-3\phi^{-1})^N.
		\label{eq:Fibonacci-rank}
	\end{align}
	An identical independent-bond product has geometric rank. For
	$r_N=Aa^N+Bb^N$,
	\begin{align}
		r_Nr_{N+2}-r_{N+1}^2
		 & =AB(ab)^N(a-b)^2.
	\end{align}
	This is nonzero in both cases. Replacing $a,b$ by $a^k,b^k$
	preserves nonvanishing for every fixed blocking length $k$.
	Unitary transformations and pure ancillas preserve rank.
\end{proof}

The scopes of these two statements are different. The rank sequences
exclude an independent bond product, not a nonproduct core tensored
with bonds. Coprime multiplicity ranks exclude a common mixed factor
only in the stated one-site sector-preserving decomposition. Blocking
changes multiplicity ranks, and circuits can redistribute sectors
between sites. Neither argument excludes every bounded-depth
spectator separation.

The categorical regrouping therefore proves that the weight-dependent
fusion spectra at the original cut can be moved into one-site
multiplicities. It does not prove that those multiplicities decouple.
This is precisely the distinction between a nontrivial canonical
fusion example and a proved example of inseparable
weight-intertwiner coupling.
